\section{Organizzazione del codice sorgente}\label{sec:organizzazione-codice}

Il repository \texttt{ecg-streaming-classifier} segue una struttura piatta, coerente con la dimensione del progetto:

\begin{verbatim}
src/main/scala/
  AppRunner.scala
  batch/EcgModelTrainer.scala
  streaming/EcgPatientSimulator.scala
  streaming/EcgStreamingPredictor.scala
  utils/Config.scala
\end{verbatim}

Il package root è \texttt{it.utiu.thesis}, definito in \texttt{build.sbt} tramite \texttt{idePackagePrefix}.
I sotto-package \texttt{batch}, \texttt{streaming} e \texttt{utils} separano responsabilità senza introdurre layer di astrazione superflui.

\subsection{Centralizzazione della configurazione}

Tutti i parametri operativi risiedono in \texttt{utils.Config}:

\begin{lstlisting}[language=Scala,caption={Parametri centralizzati in \texttt{Config.scala}.},label={lst:config}]
object Config {
  val KAFKA_BOOTSTRAP_SERVERS = "localhost:9092"
  val KAFKA_TOPIC = "ecg-input-stream"
  val SPARK_MASTER = "local[*]"

  val DATASET_PATH = "/home/andrea/data/ecg/*.csv"
  val ML_MODEL_PATH = "ml-model/ecg-pipeline-model"

  val DELTA_PREDICTIONS_PATH = "spark-warehouse/ecg_predictions"
  val STREAMING_CHECKPOINT_PATH =
    "spark-warehouse/checkpoints/ecg_predictions"
}
\end{lstlisting}

La configurazione hardcoded è un limite riconosciuto del prototipo.
In un'evoluzione produttiva, i path e gli endpoint andrebbero esternalizzati tramite Typesafe Config (HOCON), variabili d'ambiente o un servizio di configurazione centralizzato.
Il path del dataset punta a un percorso assoluto Linux (\texttt{/home/andrea/data/ecg/*.csv}): su Windows o macOS va adattato manualmente prima dell'esecuzione.

\section{Data pipeline e machine learning}\label{sec:data-pipeline-ml}

\subsection{Preprocessing e feature engineering}\label{subsec:preprocessing-fe}

Il trainer legge CSV con header e inferenza automatica dello schema.
La pulizia è minimale: \texttt{rawDF.na.drop()} elimina righe con qualsiasi valore nullo, senza imputazione.

La selezione delle colonne feature avviene per esclusione:

\begin{lstlisting}[language=Scala,caption={Selezione colonne feature in \texttt{EcgModelTrainer}.},label={lst:feature-cols}]
val featureCols = cleanDF.columns.filter(colName =>
  colName != "record" && colName != "type" && colName != "label"
)
\end{lstlisting}

Le colonne \texttt{record} (identificativo traccia) e \texttt{type} (etichetta testuale) vengono escluse dal vettore numerico.
L'esclusione di \texttt{label} è difensiva: evita leakage se la colonna fosse presente nel CSV con valori numerici già indicizzati.

Lo \texttt{StringIndexer} converte \texttt{type} in \texttt{label} numerico:

\begin{lstlisting}[language=Scala,caption={Pipeline ML: indicizzazione, assemblaggio, classificazione.},label={lst:ml-pipeline}]
val indexer = new StringIndexer()
  .setInputCol("type")
  .setOutputCol("label")
  .setHandleInvalid("keep")

val assembler = new VectorAssembler()
  .setInputCols(featureCols)
  .setOutputCol("features")
  .setHandleInvalid("keep")

val rf = new RandomForestClassifier()
  .setLabelCol("label")
  .setFeaturesCol("features")
  .setNumTrees(100)
  .setMaxDepth(15)

val pipeline = new Pipeline()
  .setStages(Array(indexer, assembler, rf))
\end{lstlisting}

\texttt{setHandleInvalid("keep")} su entrambi gli stadi è una scelta pragmatica.
In training, etichette o feature mancanti ricevono un bucket dedicato anziché causare eccezioni.
In streaming, battiti con feature incomplete possono comunque transitare nella pipeline, producendo predizioni potenzialmente inaffidabili.
Un approccio più rigoroso prevedrebbe filtraggio a monte (\texttt{drop} o quarantena su topic separato) anziché tolleranza silenziosa.

\subsection{Split, addestramento e valutazione}

\begin{lstlisting}[language=Scala,caption={Split train/test e valutazione.},label={lst:train-eval}]
val Array(trainingData, testData) =
  cleanDF.randomSplit(Array(0.8, 0.2), seed = 42L)

val pipelineModel = pipeline.fit(trainingData)

val predictions = pipelineModel.transform(testData)

val evaluator = new MulticlassClassificationEvaluator()
  .setLabelCol("label")
  .setPredictionCol("prediction")
  .setMetricName("accuracy")

val accuracy = evaluator.evaluate(predictions)
\end{lstlisting}

Lo \texttt{seed = 42L} garantisce riproducibilità dello split.
L'evaluator calcola solo l'accuratezza globale; metriche per-classe (precision, recall, F1) richiederebbero un \texttt{MulticlassClassificationEvaluator} per ogni metrica o l'uso di \texttt{BinaryClassificationMetrics} su coppie one-vs-rest.
Il Capitolo~\ref{ch:analisi-dei-risultati-e-test} estende l'analisi con metriche più granulari.

\subsection{Iperparametri Random Forest}

\texttt{numTrees=100} e \texttt{maxDepth=15} non sono stati ottimizzati con grid search.
La scelta riflette un compromesso documentato in letteratura~\cite{breiman2001random}: 100 alberi offrono varianza sufficientemente bassa senza costi eccessivi; profondità 15 limita overfitting su dataset di dimensione moderata.
Un tuning sistematico con \texttt{CrossValidator} o \texttt{TrainValidationSplit} migliorerebbe la generalizzazione, a costo di tempi di training più lunghi.

\subsection{Serializzazione e persistenza del modello}\label{subsec:serializzazione}

\begin{lstlisting}[language=Scala,caption={Salvataggio pipeline su filesystem.},label={lst:save-model}]
pipelineModel.write.overwrite().save(modelPath)
\end{lstlisting}

Spark ML salva la pipeline in una directory con struttura:

\begin{verbatim}
ml-model/ecg-pipeline-model/
  metadata/
  stages/
    0_StringIndexer_.../
    1_VectorAssembler_.../
    2_RandomForestClassificationModel_.../
\end{verbatim}

Il flag \texttt{overwrite()} consente retraining senza rimozione manuale della directory precedente.
In produzione, versionare i modelli (MLflow, DVC) eviterebbe sovrascritture accidentali e permetterebbe rollback.

\section{Gestione dello streaming}\label{sec:gestione-streaming}

\subsection{Parsing JSON e schema dinamico}\label{subsec:parsing-json}

Il consumer Kafka decodifica il payload in due passaggi:

\begin{lstlisting}[language=Scala,caption={Parsing messaggi Kafka in \texttt{EcgStreamingPredictor}.},label={lst:from-json}]
val kafkaStreamDF = spark.readStream
  .format("kafka")
  .option("kafka.bootstrap.servers", Config.KAFKA_BOOTSTRAP_SERVERS)
  .option("subscribe", Config.KAFKA_TOPIC)
  .option("startingOffsets", "earliest")
  .load()

val parsedStreamDF = kafkaStreamDF
  .selectExpr("CAST(value AS STRING) as json_string")
  .select(from_json(col("json_string"), jsonSchema).alias("data"))
  .select("data.*")
\end{lstlisting}

\texttt{startingOffsets=earliest} garantisce che, al primo avvio, vengano consumati tutti i messaggi presenti nel topic.
In restart con checkpoint, Spark riprende dall'ultimo offset committato, ignorando questa opzione.

L'alternativa professionale --- schema Avro registrato su Confluent Schema Registry --- offrirebbe validazione strict, compatibilità backward/forward e evoluzione controllata.
Per un prototipo accademico, l'inferenza schema da CSV riduce il boilerplate a costo di accoppiamento implicito.

\subsection{Trasformazione e proiezione output}

\begin{lstlisting}[language=Scala,caption={Inferenza e selezione colonne output.},label={lst:transform}]
val predictionsStreamDF = pipelineModel.transform(parsedStreamDF)

val finalStreamDF = predictionsStreamDF.select(
  col("record"),
  col("type").alias("true_label"),
  col("prediction"),
  col("probability")
)
\end{lstlisting}

La colonna \texttt{probability} contiene un vettore sparse di probabilità per classe.
In analisi successive, estrarre la probabilità della classe predetta richiede UDF o parsing del vettore.

\subsection{Checkpointing e fault tolerance}\label{subsec:checkpointing}

L'Algoritmo~\ref{alg:foreach-batch} descrive la logica di \texttt{foreachBatch}.

\begin{algorithm}[H]
  \caption{Elaborazione micro-batch con \texttt{foreachBatch}}\label{alg:foreach-batch}
  \KwIn{Micro-batch $B_i$, identificativo $batchId$, path Delta $P$, checkpoint $C$}
  \KwOut{Predizioni persistite, offset committati}
  \BlankLine
  $B_i.\text{persist}()$\;
  \If{$B_i$ non è vuoto}{
    Logga $batchId$ e mostra prime 20 righe\;
    $B_i.\text{write.format("delta").mode("append").save}(P)$\;
  }
  $B_i.\text{unpersist}()$\;
  Spark committa offset in $C$\;
\end{algorithm}

\begin{lstlisting}[language=Scala,caption={Implementazione \texttt{foreachBatch} con Delta Lake.},label={lst:foreach-batch}]
val query = finalStreamDF.writeStream
  .outputMode("append")
  .option("checkpointLocation", checkpointPath)
  .foreachBatch { (batchDF: DataFrame, batchId: Long) =>
    batchDF.persist()
    if (!batchDF.isEmpty) {
      batchDF.show(20, truncate = false)
      batchDF.write.format("delta").mode("append").save(deltaPath)
    }
    batchDF.unpersist()
    ()
  }
  .start()
\end{lstlisting}

\texttt{persist()} evita ricalcolo del batch tra \texttt{isEmpty}, \texttt{show} e \texttt{write}.
Senza persist, Spark potrebbe rieseguire la catena di trasformazioni tre volte, triplicando il costo computazionale.

\texttt{outputMode("append")} è l'unico mode compatibile con \texttt{foreachBatch} quando il sink non è nativo Spark Streaming.
Exactly-once a livello di micro-batch si ottiene dalla combinazione:
\begin{itemize}
  \item checkpoint che registra offset Kafka processati;
  \item scrittura Delta atomica per batch;
  \item restart che riprende dall'ultimo checkpoint senza riprocessare batch già committati.
\end{itemize}

\section{Storage e sink: Delta Lake}\label{sec:storage-delta}

\subsection{Configurazione sink}

I path di output sono definiti in \texttt{Config}:
\texttt{DELTA\_PREDICTIONS\_PATH} per i dati di business, \texttt{STREAMING\_CHECKPOINT\_PATH} per i metadati Spark.
La separazione è fondamentale: cancellare il checkpoint per ``resettare'' l'experimento non deve cancellare le predizioni accumulate.

\subsection{Operazioni successive}

Delta Lake supporta operazioni avanzate non implementate nel prototipo ma rilevanti per un deployment clinico:

\begin{itemize}
  \item \texttt{OPTIMIZE} --- compattazione file small per query più efficienti;
  \item \texttt{Z-ORDER BY record} --- clustering per accesso point-query;
  \item \texttt{DESCRIBE HISTORY} --- audit trail delle modifiche;
  \item \texttt{RESTORE TO VERSION} --- rollback a snapshot precedente.
\end{itemize}

Queste funzionalità sono particolarmente utili in ambito regolato, dove tracciabilità e immutabilità del log predittivo sono requisiti non negoziabili~\cite{armbrust2020delta}.

\section{Producer Kafka del simulatore}\label{sec:producer-kafka}

\subsection{Perché KafkaProducer puro}

\texttt{EcgPatientSimulator} usa il client \texttt{KafkaProducer} Java anziché il sink Spark \texttt{writeStream.format("kafka")}.
La motivazione è architetturale: in produzione, il dispositivo edge non eseguirebbe Spark.
Usare un producer standalone avvicina il simulatore a un sensore reale che pubblica autonomamente sul broker.

\begin{lstlisting}[language=Scala,caption={Configurazione e invio messaggi Kafka.},label={lst:kafka-producer}]
val props = new Properties()
props.put(ProducerConfig.BOOTSTRAP_SERVERS_CONFIG,
  Config.KAFKA_BOOTSTRAP_SERVERS)
props.put(ProducerConfig.KEY_SERIALIZER_CLASS_CONFIG,
  "org.apache.kafka.common.serialization.StringSerializer")
props.put(ProducerConfig.VALUE_SERIALIZER_CLASS_CONFIG,
  "org.apache.kafka.common.serialization.StringSerializer")

val producer = new KafkaProducer[String, String](props)

for ((jsonPayload, index) <- jsonRecords.zipWithIndex) {
  val record = new ProducerRecord[String, String](
    topic, s"beat-$index", jsonPayload)
  producer.send(record)
  Thread.sleep(1000)
}
\end{lstlisting}

\subsection{Limiti dell'implementazione}

\begin{itemize}
  \item \texttt{.collect()} materializza 1000 record sul driver: accettabile in demo, problematico in scala.
  \item Nessuna gestione errori sul \texttt{send}: callback o \texttt{get()} andrebbero aggiunti per garantire delivery.
  \item \texttt{Thread.sleep(1000)} introduce rate deterministico, non rappresentativo del jitter reale.
  \item SparkSession viene creato ma non chiuso (\texttt{spark.stop()} commentato): leak di risorse in esecuzioni ripetute.
\end{itemize}

\section{Orchestrazione e ciclo di vita}\label{sec:orchestrazione}

\texttt{AppRunner} implementa la sequenza di avvio descritta nel Capitolo~\ref{sec:app-runner}.
Il frammento saliente riguarda il thread demone del predictor:

\begin{lstlisting}[language=Scala,caption={Avvio predictor in thread demone.},label={lst:app-runner-thread}]
val predictorThread = new Thread(() => {
  EcgStreamingPredictor.main(Array.empty)
})
predictorThread.setDaemon(true)
predictorThread.start()
Thread.sleep(15000)
EcgPatientSimulator.main(Array.empty)
predictorThread.join()
\end{lstlisting}

\texttt{setDaemon(true)} fa sì che la JVM termini quando il main thread completa, anche se il predictor è ancora in \texttt{awaitAnyTermination}.
In assenza di \texttt{join()}, il programma potrebbe chiudersi prematuramente.
L'attesa di 15 secondi prima del simulatore è una euristica fragile: un health-check esplicito sarebbe preferibile.

\subsection{Banner di avvio}

\texttt{AppRunner} stampa un banner testuale all'avvio (\texttt{ECG STREAMING CLASSIFIER - BIG DATA PIPELINE}).
Funge da indicatore visivo in console, senza impatto funzionale sulla pipeline.

\section{Riepilogo implementativo}

La Tabella~\ref{tab:componenti-implementazione} collega ogni componente alle scelte implementative salienti.

\begin{table}[H]
  \centering
  \caption{Riepilogo componenti e scelte implementative.}
  \label{tab:componenti-implementazione}
  \begin{tabular}{@{}p{3.2cm}p{4.5cm}p{5cm}@{}}
    \toprule
    \textbf{Componente} & \textbf{File} & \textbf{Scelta chiave} \\
    \midrule
    Orchestratore & \texttt{AppRunner.scala} & Thread demone + sleep 15s \\
    Trainer & \texttt{EcgModelTrainer.scala} & Pipeline ML, RF 100/15, seed 42 \\
    Simulatore & \texttt{EcgPatientSimulator.scala} & KafkaProducer, 1000 msg, 1 msg/s \\
    Predictor & \texttt{EcgStreamingPredictor.scala} & from\_json, foreachBatch, Delta \\
    Config & \texttt{Config.scala} & Path hardcoded, local[*] \\
    \bottomrule
  \end{tabular}
\end{table}

Il capitolo successivo presenta i risultati sperimentali derivati da questa implementazione, con metriche di classificazione e misure di latenza e throughput.
